\documentclass{article}
\usepackage{graphicx} % Required for inserting images

\title{Sutton and Barton Exercises Chapter 10}
\author{danieltocco0 }
\date{March 2026}

\documentclass{article}
\usepackage{amsmath}
\begin{document}


\maketitle

\section{Chapter 10}
\textit{Exercise 10.1} We have not explicitly considered or given pseudocode for any Monte Carlo
methods in this chapter. What would they be like? Why is it reasonable not to give
pseudocode for them? How would they perform on the Mountain Car task?
\\
\\

The pseudocode for MC is not going to look too different than n-step SARSA.

In fact, it may be simpler, because MC will not involve the bootstrapping 

term. Consequently, because there is no bootstrapping, MC is going to 

perform poorly here. The episodes, especially early on, are long, and 

the rewards are sparse (that is, there is little, or sparse, information

when your reward signal is -1 for all actions except the one that

actually achieves the goal). N-step TD methods thus perform better since

they are constantly exploring new states (due to optimistic initial values)

and bootstrap during the actual run.
\\\\
\textit{Exercise 10.2} Give pseudocode for semi-gradient one-step Expected Sarsa for control.
\\
\\

In line with GPT and another answer I found, I updated the box from page 

244. The update applies to the weight vector, as we are taking the expected 

bootstrapped value:
\\

$\mathbf{w} \leftarrow \mathbf{w} + \alpha \left[ R + \gamma \sum_{a} \pi(a | S') \hat{q}(S', a, \mathbf{w}) - \hat{q}(S, A, \mathbf{w}) \right] \nabla \hat{q}(S, A, \mathbf{w}).$
\\

Because we still have to take an action after the update, we then take A'. 

We can flip this line and the weight vector update, if we really want to, 

though this need not necessarily be the case (we can take A' and still update 

the previous line with just R).
\\\\
\textit{Exercise 10.3} Why do the results shown in Figure 10.4 have higher standard errors at
large n than at small n?
\\
\\

Because the variance of a sum of n random variables grows with n. For large 

n then, successive runs will result in more varied learning curves.
\\\\
\textit{Exercise 10.4} Why do the results shown in Figure 10.4 have higher standard errors at
large n than at small n?
\\
\\

We just make the update to the delta term as follows: 
\\

$\delta \leftarrow R - \bar{R} + \max_{a} \hat{q}(S', a, \mathbf{w}) - \hat{q}(S, A, \mathbf{w})$.
\\

This is because Q-learning assumes you take the optimal action. To be sure,

we have an estimate of the value at State S, then action A, and then we 

capture 1) the differential reward, and 2) our prior estimate that we 

committed to and the estimate of the value assuming we act optimally after 

action A (I am basically trying to re-explain TD error, but we include 

the differential error now)!
\\\\
\textit{Exercise 10.5} What equations are needed (beyond 10.10) to specify the differential
version of TD(0)?
\\
\\

TD(0) updates the state-value functions rather than action value functions.

Similarly to 10.4 then, we make an update to the delta term as follows:
\\

$\delta = R - \bar{R} + \hat{V}(S', \mathbf{w}) - \hat{V}(S, \mathbf{w}).$

The logic is, again, similar, except we are are of course working with state 

values.
\\\\
\textit{Exercise 10.6} Suppose there is an MDP that under any policy produces the deterministic sequence of rewards +1,0,+1,0,+1,0,... going on forever. Technically, this violates
ergodicity; there is no stationary limiting distribution $\mu_{\pi}$ and the limit (10.7) does not
exist. Nevertheless, the average reward (10.6) is well defined. What is it? Now consider
two states in this MDP. From A, the reward sequence is exactly as described above,
starting with a +1, whereas, from B, the reward sequence starts with a 0 and then
continues with +1,0,+1,0,.... We would like to compute the differential values of A and
B. Unfortunately, the differential return (10.9) is not well defined when starting from
these states as the implicit limit does not exist. To repair this, one could alternatively
define the differential value of a state a
\\

$v_{\pi}(s) \doteq \lim_{\gamma \to 1} \lim_{h \to \infty} \sum_{t=0}^{h} \gamma^t \Big( \mathbb{E}_{\pi}[R_{t+1} | S_0 = s] - r(\pi) \Big)$
\\
Under this definition, what are the differential values of states A and B?
\\
\\

So we know over the long run $r(\pi) = \frac{1}{2}$. We also know B is just A shifted 

one step forward. This means the differential reward of A and B are, 

respectively: 
\\
$v_{\pi}(A) = \sum_{t=0}^{\infty} \gamma^t \Big( \mathbb{E}_{\pi}[R_{t+1} | S_0 = A] - r(\pi) \Big)$\\
$v_{\pi}(A) = \gamma^0\left(+\frac{1}{2}\right) + \gamma^1\left(-\frac{1}{2}\right) + \gamma^2\left(+\frac{1}{2}\right) + \gamma^3\left(-\frac{1}{2}\right) + \cdots$\\
$v_{\pi}(A) = \frac{1}{2}\left(1 - \gamma + \gamma^2 - \gamma^3 + \cdots\right)$\\
$v_{\pi}(A) = \frac{1}{2} \cdot \frac{1}{1+\gamma}$\\
$v_{\pi}(A) = \lim_{\gamma \to 1} \frac{1}{2} \cdot \frac{1}{1+\gamma} = \frac{1}{2} \cdot \frac{1}{2} = \frac{1}{4}$\\.
$v_{\pi}(B) = \sum_{t=0}^{\infty} \gamma^t \Big( \mathbb{E}_{\pi}[R_{t+1} | S_0 = B] - r(\pi) \Big)$\\
$v_{\pi}(B) = \gamma^0\left(-\frac{1}{2}\right) + \gamma^1\left(+\frac{1}{2}\right) + \gamma^2\left(-\frac{1}{2}\right) + \gamma^3\left(+\frac{1}{2}\right) + \cdots$\\
$v_{\pi}(B) = -\frac{1}{2}\left(1 - \gamma + \gamma^2 - \gamma^3 + \cdots\right)$\\
$v_{\pi}(B) = -\frac{1}{2} \cdot \frac{1}{1+\gamma}$\\
$v_{\pi}(B) = \lim_{\gamma \to 1} -\frac{1}{2} \cdot \frac{1}{1+\gamma} = -\frac{1}{2} \cdot \frac{1}{2} = -\frac{1}{4}$.
\\

Again, note that B is shifted one step forward from A; because, again, 

we are taking the differential value this is \textit{why} it is $-\frac{1}{4}$.
\\\\
\textit{Exercise 10.7}  Consider a Markov reward process consisting of a ring of three states A, B,
and C, with state transitions going deterministically around the ring. A reward of +1 is
received upon arrival in A and otherwise the reward is 0. What are the differential values
of the three states, using (10.13)?
\\
\\
We know $r(\pi) = \frac{1}{3}$. So we get:
\\
\begin{align*}
&\textbf{Starting at A} \text{ (next arrivals are B, C, A, B, C, A, } \ldots\text{):} \\
&\text{Rewards: } 0, 0, +1, 0, 0, +1, \ldots \\
&\text{Centered: } -\frac{1}{3}, -\frac{1}{3}, +\frac{2}{3}, -\frac{1}{3}, -\frac{1}{3}, +\frac{2}{3}, \ldots \\[10pt]
&\textbf{Starting at B} \text{ (next arrivals are C, A, B, C, A, } \ldots\text{):} \\
&\text{Rewards: } 0, +1, 0, 0, +1, 0, \ldots \\
&\text{Centered: } -\frac{1}{3}, +\frac{2}{3}, -\frac{1}{3}, -\frac{1}{3}, +\frac{2}{3}, -\frac{1}{3}, \ldots \\[10pt]
&\textbf{Starting at C} \text{ (next arrivals are A, B, C, A, } \ldots\text{):} \\
&\text{Rewards: } +1, 0, 0, +1, 0, 0, \ldots \\
&\text{Centered: } +\frac{2}{3}, -\frac{1}{3}, -\frac{1}{3}, +\frac{2}{3}, -\frac{1}{3}, -\frac{1}{3}, \ldots
\end{align*}
\\

And then actually computing the values:

\begin{align*}
&\textbf{Computing } v_{\pi}(A): \\
&v_{\pi}(A) = \left(-\frac{1}{3} - \frac{\gamma}{3} + \frac{2\gamma^2}{3}\right) \cdot \frac{1}{1-\gamma^3} \\
&\lim_{\gamma \to 1}: \text{ numerator } = -\frac{1}{3} - \frac{1}{3} + \frac{2}{3} = 0, \text{ so apply L'Hôpital's rule:} \\
&\frac{d}{d\gamma}\left(-\frac{1}{3} - \frac{\gamma}{3} + \frac{2\gamma^2}{3}\right) = -\frac{1}{3} + \frac{4\gamma}{3} \to 1 \text{ as } \gamma \to 1 \\
&\frac{d}{d\gamma}(1 - \gamma^3) = -3\gamma^2 \to -3 \text{ as } \gamma \to 1 \\
&v_{\pi}(A) = \frac{1}{-3} = -\frac{1}{3} \\[10pt]
&\textbf{Computing } v_{\pi}(B): \\
&v_{\pi}(B) = \left(-\frac{1}{3} + \frac{2\gamma}{3} - \frac{\gamma^2}{3}\right) \cdot \frac{1}{1-\gamma^3} \\
&\lim_{\gamma \to 1}: \text{ numerator } = -\frac{1}{3} + \frac{2}{3} - \frac{1}{3} = 0, \text{ so apply L'Hôpital's rule:} \\
&\frac{d}{d\gamma}\left(-\frac{1}{3} + \frac{2\gamma}{3} - \frac{\gamma^2}{3}\right) = \frac{2}{3} - \frac{2\gamma}{3} \to 0 \text{ as } \gamma \to 1 \\
&\frac{d}{d\gamma}(1 - \gamma^3) = -3\gamma^2 \to -3 \text{ as } \gamma \to 1 \\
&v_{\pi}(B) = \frac{0}{-3} = 0 \\[10pt]
&\textbf{Computing } v_{\pi}(C): \\
&v_{\pi}(C) = \left(\frac{2}{3} - \frac{\gamma}{3} - \frac{\gamma^2}{3}\right) \cdot \frac{1}{1-\gamma^3} \\
&\lim_{\gamma \to 1}: \text{ numerator } = \frac{2}{3} - \frac{1}{3} - \frac{1}{3} = 0, \text{ so apply L'Hôpital's rule:} \\
&\frac{d}{d\gamma}\left(\frac{2}{3} - \frac{\gamma}{3} - \frac{\gamma^2}{3}\right) = -\frac{1}{3} - \frac{2\gamma}{3} \to -1 \text{ as } \gamma \to 1 \\
&\frac{d}{d\gamma}(1 - \gamma^3) = -3\gamma^2 \to -3 \text{ as } \gamma \to 1 \\
&v_{\pi}(C) = \frac{-1}{-3} = \frac{1}{3}
\end{align*}
\\

Note we are using L'Hôpital's rule because if use direct substitution as 

$\lim_{\gamma \to 1}$ into our value equations, the numerator and denominator 

both evaluate to 0.
\\\\
{Exercise 10.8} The pseudocode in the box on page 251 updates $\bar{R}$ using $\delta_t$ as an error rather than simply $R_{t+1} - \bar{R}$. Both errors work, but using $\delta_t$ is better. To see why, consider the ring MRP of three states from Exercise 10.7. The estimate of the average reward should tend towards its true value of $\frac{1}{3}$. Suppose it was already there and was held stuck there. What would the sequence of $R_{t+1} - \bar{R}$ errors be? What would the sequence of $\delta_t$ errors be (using Equation 10.10)? Which error sequence would produce a more stable estimate of the average reward if the estimate were allowed to change in response to the errors? Why?
\\
\\
Consider again the differential error sequence of $v_{\pi}(A)$:
\begin{align*}
-\frac{1}{3}, -\frac{1}{3}, +\frac{2}{3}, -\frac{1}{3}, -\frac{1}{3}, +\frac{2}{3}, \ldots \\[10pt]
\end{align*}
\\

This sequence never settles. It keeps oscillating between +2/3 and -1/3

forever. If $\bar{R}$ were allowed to update in response to these errors, 

it would get yanked up and then back down every cycle — very unstable. 

Now recall that $\delta_t$ is the full TD error:
\begin{align*}
\delta_t = R_{t+1} - \bar{R} + \hat{v}(S_{t+1}, \mathbf{w}) - \hat{v}(S_t, \mathbf{w})
\end{align*}

Using the true values $v(A) = -\frac{1}{3}$, $v(B) = 0$, $v(C) = \frac{1}{3}$, and starting at C,

the transitions are $C \to A \to B \to C \to A \to \cdots$

\begin{align*}
&\textbf{Step 1: } C \to A \\
&\delta = +1 - \frac{1}{3} + v(A) - v(C) = +1 - \frac{1}{3} + \left(-\frac{1}{3}\right) - \frac{1}{3} = 0 \\[6pt]
&\textbf{Step 2: } A \to B \\
&\delta = 0 - \frac{1}{3} + v(B) - v(A) = 0 - \frac{1}{3} + 0 - \left(-\frac{1}{3}\right) = 0 \\[6pt]
&\textbf{Step 3: } B \to C \\
&\delta = 0 - \frac{1}{3} + v(C) - v(B) = 0 - \frac{1}{3} + \frac{1}{3} - 0 = 0
\end{align*}

So the sequence of $\delta_t$ errors is $0, 0, 0, 0, 0, 0, \ldots$.
\\

In short — $R_{t+1}$ - \bar{R} sees only the noisy instantaneous reward, 

while $\delta_{t}$ sees the reward plus the change in long-term prospects, 

which is a much smoother signal.
\\\\
{Exercise 10.9} In the differential semi-gradient $n$-step Sarsa algorithm, the step-size parameter on the average reward, $\beta$, needs to be quite small so that $\bar{R}$ becomes a good long-term estimate of the average reward. Unfortunately, $\bar{R}$ will then be biased by its initial value for many steps, which may make learning inefficient. Alternatively, one could use a sample average of the observed rewards for $\bar{R}$. That would initially adapt rapidly but in the long run would also adapt slowly. As the policy slowly changes, $\bar{R}$ would also change; the potential for such long-term nonstationarity makes sample-average methods ill-suited. In fact, the step-size parameter on the average reward is a perfect place to use the unbiased constant-step-size trick from Exercise 2.7. Describe the specific changes needed to the boxed algorithm for differential semi-gradient $n$-step Sarsa to use this trick.
\\
\\

We need to introduce a trace variable and then divide the step size, giving:

\begin{aligned}
o &\leftarrow o + \beta (1 - o) \\
\bar{R} &\leftarrow \bar{R} + \frac{\beta}{o}\,\delta
\end{aligned}.

As a reminder, the sample average would go like $\frac{1}{n}$. While this means

learning is fast at first, later on learning is very slow. Because the distribution 

is non-stationary because of policy improvement, we would eventually be 

making very small, if barely any, improvements towards the optimal policy. 

On the other hand, if we had a small constant step-size parameter, 

this will be good for a long-term estimate, but then $\bar{R}$ will be 

biased by its initial value for a long time because, when we unroll 

the recursion, $(1-\beta^{n})$ will decay very slowly. The unbiased 

constant-step-size trick is thus ideal because it gives us the best of both 

worlds, essentially normalizing the weights so they sum to 1 (while this 

information is in the question, I am simply repeating it with slightly different 


wording for my own benefit).




\end{document}
